\section*{Chapitre \ref{ch:in:education-pipelines} --- Parcours de formation}

\begin{solution}{exo:in:education-pipelines:1}
$4\,584/305=15.0$~: quinze fois.
\end{solution}

\begin{solution}{exo:in:education-pipelines:2}
$67/70=95.7~\%$.
\end{solution}

\begin{solution}{exo:in:education-pipelines:3}
42~: six problèmes de 7 points chacun.
\end{solution}

\begin{solution}{exo:in:education-pipelines:4}
Avec 38 non-déclarants, entre les $(49-38)/60=18.3$\up{e} et les $49/60=81.7$\up{e}~centiles de la rémunération déclarée~: la
médiane déclarée dirait peu de chose de celle de la promotion.
\end{solution}

\begin{solution}{exo:in:education-pipelines:5}
Les programmes choisissent leur code, et certains programmes d'ingénierie financière en utilisent d'autres (ingénierie,
gestion, statistiques)~: le code les laisse donc de côté~; si des programmes ont rejoint le code au fil du temps, le nombre
augmente plus vite que celui de ces diplômes.
\end{solution}

\begin{solution}{exo:in:education-pipelines:6}
$1\,760+1\,628+1\,165+454=5\,007$ doctorats de recherche en physique, informatique, mathématiques et statistiques, contre
$4\,584+870=5\,454$ masters dans les deux codes financiers~: à peu près autant.
\end{solution}

\begin{solution}{exo:in:education-pipelines:7}
Les masters d'informatique sont passés de 11\,598 à 25\,543, soit 2.20 fois~; les mathématiques financières de 3\,554 à 4\,584,
soit 1.29 fois.
\end{solution}

\begin{solution}{exo:in:education-pipelines:8}
Les deux rapports utilisent des dénominateurs et des normes différents~: l'un ne dit pas combien de rémunérations ont été
déclarées, l'autre compte l'emploi de courte durée comme insertion selon sa norme~; les promotions sortent à des moments
différents de l'année et sur des marchés différents~; et une médiane du \omterm{def:in:how-pay-works:base}{salaire de base} ne dit rien des primes ni de la
nature des emplois. Un écart de 5\,000~\$ peut tenir à ces choix.
\end{solution}

\begin{solution}{pb:in:education-pipelines:1}
\begin{enumerate}
\item Comme dans la définition du chapitre.
\item 27.0305 mathématiques financières («~l'application des mathématiques et de la statistique à l'industrie financière~»), et
30.7104 analytique financière («~modélisation des mégadonnées financières~»), nouveau en 2020.
\item Diplômes décernés par établissement, code, spécialité et niveau, de juillet à juin.
\item Pour qu'un étudiant ayant deux spécialités ne soit compté qu'une fois.
\item D'autres programmes utilisent d'autres codes.
\item 305, 773, 3\,554 et 4\,584.
\item Comme dans l'encadré des diplômes.
\item Moins du cinquième des masters d'informatique.
\item 1\,760, 1\,628, 1\,165 et 454~: 5\,007.
\item Recherche issue des doctorats et des masters généralistes~; ingénierie issue de l'informatique~; quant bancaire et
structuration issus également du master spécialisé.
\item Comme dans la définition du chapitre.
\item 70 demandes, 67 acceptées, 150\,000~\$, nombre de déclarants non précisé~; 98 demandes, 98 acceptées dans les trois
mois (emploi de courte durée compris), 89 déclarants, 145\,000~\$.
\item Entre les quantiles déclarés $(n/2-(n-r))/r$ et $(n/2)/r$~: 44.9~\% et 55.1~\% pour 89 sur 98.
\item Ce que gagnaient les non-déclarants et pourquoi ils n'ont pas déclaré~; les primes~; ce qui compte comme un emploi.
\item Rapidité et profondeur sur des problèmes bien posés~; pas la recherche sur des problèmes mal posés.
\item 5\,454 masters en 2024 contre 1\,466 dossiers de visa pour l'exercice 2025~: 3.72.
\item Les deux décomptes portent sur des personnes différentes~: tous les diplômés, toutes nationalités, contre les dossiers
de visa déposés par un échantillon d'employeurs.
\item Diplômés en informatique, statistiques, mathématiques et physique, ainsi que les autres programmes financiers.
\item Ses dénominateurs (taille de la promotion, demandes, acceptations, déclarants), sa norme de déclaration et les dates.
\item Le marché compte plusieurs milliers de diplômés spécialisés par an, à côté de dizaines de milliers de diplômés généralistes, et
les entreprises recrutent dans les deux populations. Un nouveau programme sera jugé sur les preuves d'insertion qu'il peut présenter avec ses dénominateurs.
\end{enumerate}
\end{solution}

\begin{solution}{iq:in:education-pipelines:1}
Un exemple précis de travail sur un problème sans réponse connue pendant des mois~: impasses, vérification de ses propres
résultats, et choix de ce à quoi renoncer.
\iqlookfor{preuves de recherche indépendante.}
\end{solution}

\begin{solution}{iq:in:education-pipelines:2}
Présenter un projet qui peut être relancé, et expliquer un choix de conception, ce qui a été essayé et ce qui a été mesuré pour trancher.
\iqlookfor{raisonnement sur les compromis, pas une liste d'outils.}
\end{solution}

\begin{solution}{iq:in:education-pipelines:3}
Sur combien, combien ont cherché un emploi, ce qui compte comme une insertion (emploi de courte durée~?), à quelle date, et
combien ont déclaré leur rémunération.
\iqlookfor{dénominateurs et définitions.}
\end{solution}

\begin{solution}{iq:in:education-pipelines:4}
Une estimation de Fermi par l'offre (diplômés des filières d'alimentation multipliés par la part embauchée) et par la
demande (effectifs des entreprises multipliés par le turnover et la croissance), les deux étant comparées~; préciser l'incertitude.
\iqlookfor{deux estimations indépendantes.}
\end{solution}

\begin{solution}{iq:in:education-pipelines:5}
Avec un code que d'autres peuvent lire et exécuter~: un dépôt avec des tests, une construction, de la documentation et un historique, et une
partie qu'on peut expliquer ligne par ligne.
\iqlookfor{des preuves plutôt que des affirmations.}
\end{solution}

\begin{solution}{iq:in:education-pipelines:6}
Chercher une baisse correspondante dans les codes qu'il a remplacés, comparer les établissements qui ont déclaré sous les deux codes, et vérifier des totaux qui
ne dépendent pas du code.
\iqlookfor{passerelles entre codes et totaux.}
\end{solution}
